
\documentclass[border=4pt]{standalone}
\usepackage{polyglossia}
\setdefaultlanguage{russian}
\setotherlanguage{english}
\usepackage{fontspec}
\IfFontExistsTF{CMU Serif}{\setmainfont{CMU Serif}}{\setmainfont{Liberation Serif}}
\IfFontExistsTF{CMU Sans Serif}{\setsansfont{CMU Sans Serif}}{\setsansfont{Liberation Sans}}
\IfFontExistsTF{CMU Typewriter Text}{\setmonofont{CMU Typewriter Text}}{\setmonofont{DejaVu Sans Mono}}
\usepackage{amsmath,amssymb,mathtools,bm}
\usepackage{xcolor}
\definecolor{accent}{HTML}{1F4E79}
\definecolor{accent2}{HTML}{B85450}
\definecolor{shade}{HTML}{F4F1EA}
\colorlet{prosvBlue}{accent}
\colorlet{prosvRed}{accent2}
\colorlet{prosvLight}{accent!12}
\colorlet{prosvCream}{shade}
\definecolor{prosvGold}{HTML}{E8B647}
\definecolor{prosvGreen}{HTML}{6A9F58}
\colorlet{prosvGray}{black!55}
\usepackage{tikz}
\usetikzlibrary{calc, arrows.meta, decorations.pathreplacing,
                positioning, shapes.geometric, shapes.misc, fit,
                backgrounds}
\usepackage{pgfplots}
\pgfplotsset{compat=1.18}
\newcommand{\R}{\mathbb{R}}
\newcommand{\N}{\mathbb{N}}
\newcommand{\Z}{\mathbb{Z}}
\newcommand{\eps}{\varepsilon}
\newcommand{\norm}[1]{\left\|#1\right\|}
\newcommand{\abs}[1]{\left|#1\right|}
\newcommand{\NOD}{\gcd}
\newcommand{\Eul}{\varphi}
\newcommand{\Zn}[1]{\mathbb{Z}_{#1}}
\newcommand{\modn}[1]{\,(\mathrm{mod}\,#1)}
\newcommand{\term}[1]{\textcolor{accent2}{\textbf{#1}}}
\newcommand{\engterm}[1]{\textit{#1}}
\DeclareMathOperator*{\argmin}{arg\,min}
\DeclareMathOperator*{\argmax}{arg\,max}
\begin{document}

\begin{tikzpicture}[
  >=Stealth, every node/.style={font=\footnotesize},
  uni/.style={draw=prosvBlue, fill=prosvLight, rounded corners, minimum height=0.5cm, minimum width=2.2cm},
  cell/.style={draw=prosvRed!70, fill=prosvCream, minimum height=0.6cm, minimum width=0.7cm},
]
  \node[uni] (k1) at (0, 2)   {ivan@mail};
  \node[uni] (k2) at (0, 1.3) {olga@mail};
  \node[uni] (k3) at (0, 0.6) {petr@yand};
  \node[uni] (k4) at (0, -0.1){anna@bk};
  \node[uni] (k5) at (0, -0.8){vova@mail};

  \node[font=\footnotesize, prosvBlue] at (0, 2.7) {ключи $x$ (вселенная $U$)};

  \foreach \i in {0,1,2,3,4,5,6,7} {
    \node[cell] (c\i) at (5 + 0.75*\i, 0.5) {\i};
  }
  \node[font=\footnotesize, prosvBlue] at (7.5, 1.3) {хеш-таблица $T$ (размер $m=8$)};

  \draw[->, prosvGreen!60!black] (k1.east) -- (c3.west);
  \draw[->, prosvGreen!60!black] (k2.east) -- (c0.west);
  \draw[->, prosvGreen!60!black] (k3.east) -- (c5.west);
  \draw[->, prosvGreen!60!black] (k4.east) -- (c1.west);
  \draw[->, prosvRed!80, very thick] (k5.east) -- (c3.west);

  \node[font=\footnotesize, prosvRed, align=center] at (4, -1.2) {две стрелки в одну\\ ячейку $=$ коллизия};
\end{tikzpicture}
\end{document}
